% ============================================================
%  Chapter 5: Design
%  Included from main.tex via % ============================================================
%  Chapter 5: Design
%  Included from main.tex via % ============================================================
%  Chapter 5: Design
%  Included from main.tex via % ============================================================
%  Chapter 5: Design
%  Included from main.tex via \input{Chapters/Design}.
%  Do NOT add \documentclass or \begin{document} here.
%  Follows Chapter 4 (Specification); precedes Chapter 6 (Implementation).
%  Scope: architecture, design decisions with named alternatives and
%  rejection reasons, pipeline and component design, data model,
%  normalisation, evaluation framework, execution model, maintenance,
%  traceability. Written before implementation; "expected deviations"
%  section flags constraints that will shape Chapter 6.
% ============================================================

\chapter{Design}
\label{chap:design}

% ------------------------------------------------------------
\section{From specification to design}
\label{sec:design-bridge}

Chapter~\ref{chap:spec} fixed what the system must satisfy;
this chapter states how it is realised and what was rejected on the
way. Every design decision below is anchored in a requirement from
\S\ref{sec:spec-fr}--\S\ref{sec:spec-experimental-requirements} and
traces upstream to a gap named in \S\ref{sec:lr-gap}. Four
requirements shape the chapter. ER1 (modality isolation) dictates
that the two pipelines share every component except the one under
test, which constrains the architectural boundary in \S\ref{sec:design-overview}. ER2 (multilingual and prompt-controlled
coverage) dictates that the prompt and the LLM are auditable
artefacts, not inlined strings, which forces the component
abstractions in \S\ref{sec:design-components}. ER6 (null-aware
evaluation) dictates that both predictions and ground truth are
normalised identically before comparison, which drives the
normalisation design in \S\ref{sec:design-normalisation}. ER7
(cross-provider generality) dictates that swapping a provider cannot
change pipeline code, which drives the registry in
\S\ref{sec:design-components}. The remaining decisions fall out from
these four.

% ------------------------------------------------------------
\section{System overview}
\label{sec:design-overview}

Figure~\ref{fig:design-arch} gives the architecture. Three
boundaries matter. \textit{Shared input}
(manifest, loader, PDF, rasteriser) is identical across every
condition, so any difference in raw pixel content is a data-source
effect rather than a pipeline effect. \textit{Modality branch}
contains the single step that varies: either Tesseract followed by
the LLM in text mode, or the LLM in vision mode. \textit{Shared
output} (JSON parse, normalise, validate, persist, evaluate) is
identical across every condition, so any difference in final output
distribution is attributable to the branch content rather than to
the finalisation logic.

The branch is the only place ER1 can be violated, which is why it
is drawn as a single box rather than as two parallel pipelines: two
parallel pipelines have two places to drift. Holding the boundary
tight enough for one provider to be swapped independently of one
modality (\S\ref{sec:design-components}) costs some visual symmetry
in the figure but removes the dominant confound identified in
\S\ref{sec:lr-controlled}.

\begin{figure}[ht]
    \centering
    \includegraphics[width=\linewidth]{Chapters/figures/fig-5-1-architecture.png}
    \caption{End-to-end architecture. Shaded regions are held
    constant across conditions; the central dashed region is the
    modality branch where the input differs. The design isolates
    the ER1 control to a single step rather than to a whole
    pipeline.}
    \label{fig:design-arch}
\end{figure}

Figure~\ref{fig:design-seq} refines the same picture at the
per-document scale. The two paths differ only in whether Tesseract
sits between the rasteriser and the LLM; every call on either side
of that step is identical. The timing keys in brackets match the
per-stage measurements written to \texttt{timings.json} and drive
the latency analysis for RQ2.

\begin{figure}[ht]
    \centering
    \includegraphics[width=\linewidth]{Chapters/figures/fig-5-2-sequence-diagram.png}
    \caption{Per-document sequence for both modalities. The only
    step that varies is whether Tesseract sits in the chain.}
    \label{fig:design-seq}
\end{figure}

% ------------------------------------------------------------
\section{Design decisions and alternatives considered}
\label{sec:design-decisions}

Table~\ref{fig:design-decisions} summarises the ten decisions that
realise the specification, names the alternative considered for
each, and records the specific technical reason for rejection. The
rest of this chapter expands on the decisions that drive the
experiment outcome; the remaining rows are covered by the table.

\begin{figure}[ht]
    \centering
    \includegraphics[width=\linewidth]{Chapters/figures/fig-5-4-decisions-table.png}
    \caption{Design decisions, rejected alternatives and
    traceability to FR, NFR and ER classes.}
    \label{fig:design-decisions}
\end{figure}

\paragraph{Modality-isolation boundary.} The single branch point
keeps ER1 mechanical rather than editorial. Two parallel pipelines
would satisfy the letter of ER1 on day one and violate it silently
on any later change to, say, the prompt, the post-processor or the
JSON parser. The chosen shape makes provenance of any difference
directly inspectable: if the shared input and shared output are the
same byte-for-byte, the difference must come from the branch.

\paragraph{Cross-provider symmetry.} ER7 requires that swapping a
provider cannot alter pipeline code. The \texttt{LLMProvider} ABC
plus registry in \S\ref{sec:design-components} satisfies this
structurally: the pipeline sees only the ABC, so every provider is
exchanged at the level of a config key. The alternative was a
single adapter class with \texttt{if provider == "openai"}
branches; it was rejected because every provider-specific test
would then need to touch pipeline logic, and a future third
provider would require editing the pipeline rather than adding a
file.

\paragraph{Null-aware classification.} ER6 requires the four
classes of the taxonomy defined in \S\ref{sec:spec-evaluation}. The
classification does not live only in the evaluator: it shapes the
failure-diagnosis heuristic in \S\ref{sec:design-evalframework} and
the reporting in Chapter~\ref{chap:evaluation}. The rejected
alternative was a binary exact-match with both-null counted as
correct; it was rejected because on sparse bills (where many
fields are legitimately absent) it can return above 80\% accuracy
on a model that outputs \texttt{null} for every field, which is the
opposite of the discrimination the RQs require.

\paragraph{JSON parse hardening.} Providers do not agree on output
format (\S\ref{sec:lr-eval}). OpenAI Structured Outputs guarantees
a schema-valid JSON object. Anthropic without Structured Outputs
returns text that may be wrapped in markdown fences and may contain
preamble. A retry-with-repair prompt would satisfy this trivially
by asking the model to reformat, but at the cost of changing the
prompt between conditions and breaking ER2. The chosen design is
a four-stage offline fallback (\S\ref{sec:design-execution}) that
never reissues a call.

% ------------------------------------------------------------
\section{Pipeline architecture}
\label{sec:design-pipeline}

\subsection{Shared preprocessing}
\label{sec:design-pipeline-preproc}

Every document enters a single rasteriser before the branch, so the
pixel source for both paths is bit-identical. The alternative would
be to rasterise once per path; it was rejected because PDF
rasterisation is non-deterministic across invocations in the
general case (font hinting, anti-aliasing) and two independent
rasterisations would introduce a second unintended variable.

\subsection{Modality branch}

The OCR path runs Tesseract with the language taken from the
manifest; only the recognised text is forwarded to the LLM. The
vision path sends the first two rasterised pages directly to the
LLM in vision mode. The two-page cap follows from the scope
commitment in \S\ref{sec:spec-scope}: most target fields appear on
the first page or page two. The rejected alternative was to pass
every page; it was rejected because vision tokens scale
near-linearly with page area and the cost would grow quadratically
with document size for no demonstrated accuracy benefit on target
fields.

\subsection{Shared post-processing}

Every LLM response passes through the same four-stage chain: JSON
parse, field-level normalisation, Pydantic validation, per-document
persistence. Parsing and validation are separated so that a parse
failure is distinguishable from a schema violation: the former
keeps the raw output and retries offline
(\S\ref{sec:design-execution}), the latter fails loudly at
evaluation time. Pydantic is chosen over raw \texttt{jsonschema}
because the same model serves as the OpenAI Structured Outputs
schema, the LLM-output container, the normalisation target, the
evaluator input and the reporting envelope; a single Python class
removes five format-translation layers that would otherwise need to
be kept in sync.

% ------------------------------------------------------------
\section{Component design}
\label{sec:design-components}

Figure~\ref{fig:design-strategy} shows the two strategy hierarchies
and the registry that binds them. Every component visible to the
pipeline is an ABC; every concrete implementation is opaque to the
pipeline and instantiated at run time by
\texttt{registry.get\_provider(config)} (lazy import).

\begin{figure}[ht]
    \centering
    \includegraphics[width=\linewidth]{Chapters/figures/fig-5-3-strategy-pattern.png}
    \caption{Provider and OCR strategy hierarchies bound at
    run time by the registry. The rejected alternative
    (hard-coded dispatch) is shown beneath.}
    \label{fig:design-strategy}
\end{figure}

\subsection{LLM provider abstraction}
\label{sec:design-components-llm}

\texttt{LLMProvider} exposes \texttt{extract\_from\_text},
\texttt{extract\_from\_image} and \texttt{get\_model\_id}. Both
extraction methods return a uniform triple
(\texttt{raw\_output}, \texttt{token\_usage}, \texttt{latency\_ms}).
The pipeline calls only the ABC; provider identity enters through
the registry and leaves in the persisted \texttt{model\_id} so that
downstream evaluation can slice by model. The rejected alternative
was a hard-coded \texttt{if}/\texttt{elif} dispatch inside
\texttt{pipeline.py}; its rejection reason is recorded in
Figure~\ref{fig:design-decisions}.

\subsection{OCR engine abstraction}
\label{sec:design-components-ocr}

\texttt{OCREngine} exposes \texttt{extract\_text(image, lang)} and
is implemented by \texttt{TesseractOCR}. Tesseract is the chosen
concrete engine because it executes locally (no confound with a
second paid API in the OCR path) and ships language packs for the
four languages in scope. The abstraction survives the choice: a
later experiment can substitute EasyOCR or PaddleOCR by adding a
module without touching the pipeline, which is what ER5's
failure-attribution story needs to remain falsifiable.

\subsection{Prompt system}
\label{sec:design-components-prompt}

Prompts are versioned files in \texttt{prompts/}, selected by
\texttt{llm.prompt\_file}. The schema description is not written
into the prompt; it is generated from
\texttt{BillExtraction.schema\_description()} and injected through
the \texttt{\{schema\_description\}} placeholder at render time.
The rendered prompt is archived per run (\S\ref{sec:design-execution}) so a diff-only audit proves prompt
identity across conditions, which is what ER2 requires. Inline
string constants were rejected because differences between
conditions cannot be audited by reading a file.

\subsection{Provider-API asymmetry}
\label{sec:design-components-api}

The two providers use different APIs and different mechanisms for
structured output, and the design records the asymmetry rather
than papering over it. OpenAI is called through the Responses API
with Structured Outputs: \texttt{text\_format=BillExtraction}
constrains decoding to schema-valid JSON. Anthropic is called
through the Messages API without Structured Outputs; the same
prompt instructs a JSON response, and the provider's tendency to
wrap JSON in markdown fences is handled by
\texttt{\_strip\_json\_fencing()} before the shared parser runs.

This asymmetry is deliberate. Forcing both providers into
prompt-instructed JSON for the sake of parity would discard the
capability that gives OpenAI guaranteed schema adherence and shift
the risk profile of the experiment silently. The chosen design
preserves each provider's strongest available output mode and
records the raw LLM response per document (NFR8) so that format
failure is separable from content failure in post-hoc analysis. The
evaluation chapter revisits this asymmetry explicitly; a design
choice that hides it would make that analysis impossible.

\subsection{Configuration and CLI}

The runtime is configured by a YAML file (\texttt{config/default.yaml}) with CLI deep-merge overrides
for the handful of flags exercised per run (\texttt{--mode},
\texttt{--provider}, \texttt{--model}, \texttt{--force},
\texttt{--manifest}). The effective configuration after merging is
archived with the run. An environment-variable-only alternative was
rejected because at run-time audit it lacks the single-file
self-description that reproducibility (ER4) assumes.

% ------------------------------------------------------------
\section{Data model}
\label{sec:design-datamodel}

One Pydantic class, \texttt{BillExtraction}
(Table~\ref{fig:schema}), is the contract between prompt, LLM
output, normaliser, evaluator and reporter. Every field is
\texttt{Optional}, which is not a defensive default but the shape
the 4-class null taxonomy assumes: a distinct \texttt{null} is the
only way \texttt{hallucination} and \texttt{omission} can be
separated from \texttt{incorrect}. \texttt{PipelineResult} bundles
the extraction with the artefacts the evaluator and the
cost/latency accounting need (raw LLM output, OCR text, timings,
token usage, model id, pipeline mode); it is not a database row.
\texttt{DocumentEntry} is the resolved manifest row with its file
paths already verified; only entries that pass the six-step manifest
validation (\S\ref{sec:spec-experimental}) reach the pipeline.

Per-layer dataclasses were rejected because every field rename
would then be a five-file change (prompt description, LLM output
parser, normaliser dispatch, evaluator comparison, report column
header) with silent divergence if one file is missed; a single
Pydantic class fails the rename loudly.

% ------------------------------------------------------------
\section{Normalisation}
\label{sec:design-normalisation}

The same normaliser runs over predictions and ground truth before
any comparison. This is the one design choice that most directly
protects the meaning of ``correct'' against surface variation:
\texttt{"\texteuro"} and \texttt{"EUR"}, \texttt{"01/03/2024"} and
\texttt{"2024-03-01"}, \texttt{"OVO Energy Ltd"} and \texttt{"OVO
Energy"} all collapse to the same value before they are compared.
Normalising predictions alone was rejected because it would push
format parity onto the annotators: the set of correct surface forms
would be whatever the normaliser happened to recognise, not the set
of meaningfully correct forms. That trade breaks ER6.

Normalisation dispatches through \texttt{FIELD\_NORMALISERS}, a
mapping from field name to normaliser function. A monolithic
\texttt{normalise(fields)} with a chain of \texttt{if} statements
was rejected because adding a new field would then require editing
the normaliser, the comparator and the reporter in three places;
the dispatch table localises the addition to one entry in one table.
The float parser is the one non-trivial normaliser: it distinguishes
EU and US decimal conventions from the positions of the group and
decimal separators, falling back to standard parsing when both are
absent. LLM-as-judge comparison was rejected as an alternative
because the evaluator becomes non-reproducible the moment it calls
a remote model, breaking ER4.

% ------------------------------------------------------------
\section{Evaluation framework}
\label{sec:design-evalframework}

The evaluation is specified declaratively in
\S\ref{sec:spec-evaluation}; this section states its decision-tree
realisation.

\subsection{Metric set}

Field-level accuracy (over documents), document-level accuracy
(fraction of documents with no incorrect, hallucinated or omitted
field), overall accuracy (all fields across all documents), and
slice accuracy (by language, by utility type) are the four reported
metrics. The split addresses the failure mode named in
\S\ref{sec:lr-eval}: a single aggregate number can hide systematic
language-specific degradation, and coarse-grained metrics that
aggregate across a 12-field schema can hide field-class-specific
degradation (for example, a pipeline that is accurate on
\texttt{provider\_name} but randomly wrong on \texttt{due\_date}).
The six canonical charts listed in \S\ref{sec:spec-evaluation} are
each pinned to a specific research question. Overall-accuracy and
per-field-accuracy answer RQ1 (modality). Per-stage timing and
cost-per-document answer RQ2 (cost and latency).
Accuracy-by-language and accuracy-by-utility type answer RQ3
(multilingual and cross-issuer degradation). No chart exists
without a question.

\subsection{Null-aware classification}

Each field-level comparison emits one of five classifications
(Table~\ref{fig:null-taxonomy}): \texttt{correct},
\texttt{incorrect}, \texttt{hallucination}, \texttt{omission},
\texttt{both-null}. The first four are counted in the accuracy
denominator; \texttt{both-null} is counted as a match but excluded
from the denominator. The exclusion is what prevents a
\texttt{null}-emitting pipeline from scoring highly on sparse
bills. The classification is decided at compare time and persisted
so that the failure-diagnosis heuristic below can operate on it
without re-running comparison logic.

\subsection{Failure diagnosis}
\label{sec:design-evalframework-diagnosis}

For each non-correct field in an OCR-mediated condition, the
heuristic reads the stored \texttt{ocr\_text.txt} and searches it,
case-insensitively, for the ground-truth value. A hit means
Tesseract captured the value and the LLM misused it: the failure
is attributed to the LLM reasoning step. A miss means Tesseract did
not capture the value and no downstream stage could have succeeded:
the failure is attributed to OCR. In vision-mediated conditions no
intermediate text artefact exists, so all non-correct fields are
attributed to the vision model; the asymmetry is recorded and
discussed in Chapter~\ref{chap:evaluation} rather than hidden behind
a weaker heuristic.

Training a classifier for failure attribution was rejected because
it makes the evaluator itself an unvalidated ML component:
disagreements between the classifier and ground truth would then
need their own taxonomy, and an error in the attribution classifier
would be indistinguishable from an error in the pipeline under
test. The substring heuristic is inspectable; its known
false-positive modes (repeated digits on dates) are recorded in
Chapter~\ref{chap:evaluation}.

\subsection{Cross-run comparator}

\texttt{compare\_runs()} reads per-run \texttt{evaluation.json}
files and assembles the accuracy matrix, performance matrix, null
analysis, and slice comparisons for every run under comparison.
The output feeds the factorial analysis specified in
\S\ref{sec:spec-experimental}: the four cells of the $2 \times 2$
matrix are four runs to the comparator, and the comparator is the
only place accuracy across runs is aggregated. No comparison is
ever computed inside the pipeline, so pipeline changes cannot
accidentally rewrite comparison logic.

% ------------------------------------------------------------
\section{Execution model and reproducibility}
\label{sec:design-execution}

Figure~\ref{fig:design-rundir} gives the per-run directory layout.
A run produces one directory, named
\texttt{\{timestamp\}\_\{provider\}\_\{model\}\_\{mode\}}, that
contains the effective configuration, the manifest snapshot at run
time, the exact prompt template, a summary, and a
\texttt{documents/} subdirectory with one folder per document.
Every document folder holds the normalised extraction, the raw LLM
response, the OCR text when applicable, per-stage timings, token
and cost metadata, and on failure an \texttt{error.json} with the
full traceback. The evaluator is a consumer of this directory: it
reads only from disk, never from in-memory objects owned by the
pipeline, which is why pipeline and evaluator are independently
runnable.

\begin{figure}[ht]
    \centering
    \includegraphics[width=0.9\linewidth]{Chapters/figures/fig-5-6-run-directory.png}
    \caption{Per-run directory layout. The evaluator reads only
    from this directory; pipeline and evaluator never share memory,
    so each can be run, replayed and inspected independently.}
    \label{fig:design-rundir}
\end{figure}

Error isolation is mechanical rather than editorial: each document
is processed inside a \texttt{try}/\texttt{except}, and on failure
an \texttt{error.json} with full traceback is written and the run
proceeds. The evaluator skips any document with an
\texttt{error.json}. A crash on document forty therefore loses one
document rather than forty-one.

Figure~\ref{fig:design-parse} gives the JSON parse hardening chain.
The chain is local and offline: each stage rewrites the input
locally and retries the parse, and the final stage raises with a
diagnostic excerpt of the raw output (already archived for the
run). The rejected alternative was a retry-with-repair round-trip
to the model; it was rejected because it mutates the prompt
conditional on provider behaviour, and any prompt mutation breaks
ER2.

\begin{figure}[ht]
    \centering
    \includegraphics[width=0.75\linewidth]{Chapters/figures/fig-5-5-json-parse-fallback.png}
    \caption{JSON parse hardening chain. Every stage is local and
    offline; no stage reissues a call to the model.}
    \label{fig:design-parse}
\end{figure}

Resume is the default: if a document already has an
\texttt{extraction.json}, it is skipped. The \texttt{--force} flag
re-runs everything. This makes partial re-runs cheap, which matters
because API cost per condition is non-trivial (\S\ref{sec:spec-experimental}).

% ------------------------------------------------------------
\section{Maintenance considerations}
\label{sec:design-maintenance}

Five dependency risks are named rather than implied. Poppler
(binary dependency of \texttt{pdf2image}) and Tesseract (with the
four language packs in scope) are system-level dependencies that
must be present on every run host; the \texttt{README} and
\texttt{requirements.txt} list both. Provider SDKs
(\texttt{openai}, \texttt{anthropic}) are pinned in
\texttt{requirements.txt} because both providers make incompatible
API changes between major versions. Pydantic is pinned to v2
because OpenAI Structured Outputs uses v2-specific schema
metadata. The chart backend (\texttt{matplotlib}) is pinned loosely
because its failure mode is visual, not semantic.

Three extensibility points cover the expected directions of
future change. Adding a provider requires one new module and one
entry in \texttt{registry.\_PROVIDERS}. Adding an OCR engine
requires one new \texttt{OCREngine} subclass. Adding a schema field
requires one Pydantic attribute plus one entry in
\texttt{FIELD\_NORMALISERS}; the evaluator and reporter iterate
over \texttt{BillExtraction.SCHEMA\_FIELDS}, so they pick up new
fields automatically.

\paragraph{Expected deviations.} Three deviations from this design
are anticipated and will be recorded in Chapter~\ref{chap:impl}~\S\ref{sec:impl-deviations}.
First, the provider-API asymmetry in
\S\ref{sec:design-components-api} is itself a deviation from the
na{\"i}ve design of ``one call shape per provider''; the deviation
is documented here because hiding it would invalidate the
evaluation. Second, some providers may not return
\texttt{latency\_ms} uniformly across SDK versions, and the Timer
wrap in \texttt{src/performance.py} is expected to supersede the
provider field. Third, Tesseract language packs are not
universally installed on arbitrary hosts, and a missing pack
downgrades an OCR run to empty text; this is a ground-truth issue
for that condition, not a bug, and will be handled by the
failure-attribution heuristic.

% ------------------------------------------------------------
\section{Traceability summary}
\label{sec:design-traceability}

Figure~\ref{fig:design-trace} gives the end-to-end chain for each
experimental-requirement class. Every row connects a literature
gap to a requirement, to the design realisation in this chapter,
to the test or evaluation hook that will exercise it. The table is
the single defensible artefact for Trait~5: a marker can verify
the chain without reading the chapters in sequence.

\begin{figure}[ht]
    \centering
    \includegraphics[width=\linewidth]{Chapters/figures/fig-5-7-traceability-matrix.png}
    \caption{End-to-end traceability from literature gap to
    evaluation hook for every experimental requirement.}
    \label{fig:design-trace}
\end{figure}

% ------------------------------------------------------------
\section{Summary and forward bridge}
\label{sec:design-summary}

This chapter turned each requirement of Chapter~\ref{chap:spec}
into a design decision, named an alternative for each decision, and
stated the technical reason for rejection. The single architectural
commitment that dominates the chapter is the modality-isolation
boundary: every other decision exists to make that boundary
auditable by someone other than the author.
Chapter~\ref{chap:impl} takes this design and builds it, recording
the real deviations encountered against the expected deviations in
\S\ref{sec:design-maintenance}. Chapter~\ref{chap:evaluation} tests
whether the design produces the results predicted by the
hypotheses in \S\ref{sec:spec-objectives}, including the
alternative-design comparison that returns to the rejected rows of
Figure~\ref{fig:design-decisions} and evaluates them against the
same data.
.
%  Do NOT add \documentclass or \begin{document} here.
%  Follows Chapter 4 (Specification); precedes Chapter 6 (Implementation).
%  Scope: architecture, design decisions with named alternatives and
%  rejection reasons, pipeline and component design, data model,
%  normalisation, evaluation framework, execution model, maintenance,
%  traceability. Written before implementation; "expected deviations"
%  section flags constraints that will shape Chapter 6.
% ============================================================

\chapter{Design}
\label{chap:design}

% ------------------------------------------------------------
\section{From specification to design}
\label{sec:design-bridge}

Chapter~\ref{chap:spec} fixed what the system must satisfy;
this chapter states how it is realised and what was rejected on the
way. Every design decision below is anchored in a requirement from
\S\ref{sec:spec-fr}--\S\ref{sec:spec-experimental-requirements} and
traces upstream to a gap named in \S\ref{sec:lr-gap}. Four
requirements shape the chapter. ER1 (modality isolation) dictates
that the two pipelines share every component except the one under
test, which constrains the architectural boundary in \S\ref{sec:design-overview}. ER2 (multilingual and prompt-controlled
coverage) dictates that the prompt and the LLM are auditable
artefacts, not inlined strings, which forces the component
abstractions in \S\ref{sec:design-components}. ER6 (null-aware
evaluation) dictates that both predictions and ground truth are
normalised identically before comparison, which drives the
normalisation design in \S\ref{sec:design-normalisation}. ER7
(cross-provider generality) dictates that swapping a provider cannot
change pipeline code, which drives the registry in
\S\ref{sec:design-components}. The remaining decisions fall out from
these four.

% ------------------------------------------------------------
\section{System overview}
\label{sec:design-overview}

Figure~\ref{fig:design-arch} gives the architecture. Three
boundaries matter. \textit{Shared input}
(manifest, loader, PDF, rasteriser) is identical across every
condition, so any difference in raw pixel content is a data-source
effect rather than a pipeline effect. \textit{Modality branch}
contains the single step that varies: either Tesseract followed by
the LLM in text mode, or the LLM in vision mode. \textit{Shared
output} (JSON parse, normalise, validate, persist, evaluate) is
identical across every condition, so any difference in final output
distribution is attributable to the branch content rather than to
the finalisation logic.

The branch is the only place ER1 can be violated, which is why it
is drawn as a single box rather than as two parallel pipelines: two
parallel pipelines have two places to drift. Holding the boundary
tight enough for one provider to be swapped independently of one
modality (\S\ref{sec:design-components}) costs some visual symmetry
in the figure but removes the dominant confound identified in
\S\ref{sec:lr-controlled}.

\begin{figure}[ht]
    \centering
    \includegraphics[width=\linewidth]{Chapters/figures/fig-5-1-architecture.png}
    \caption{End-to-end architecture. Shaded regions are held
    constant across conditions; the central dashed region is the
    modality branch where the input differs. The design isolates
    the ER1 control to a single step rather than to a whole
    pipeline.}
    \label{fig:design-arch}
\end{figure}

Figure~\ref{fig:design-seq} refines the same picture at the
per-document scale. The two paths differ only in whether Tesseract
sits between the rasteriser and the LLM; every call on either side
of that step is identical. The timing keys in brackets match the
per-stage measurements written to \texttt{timings.json} and drive
the latency analysis for RQ2.

\begin{figure}[ht]
    \centering
    \includegraphics[width=\linewidth]{Chapters/figures/fig-5-2-sequence-diagram.png}
    \caption{Per-document sequence for both modalities. The only
    step that varies is whether Tesseract sits in the chain.}
    \label{fig:design-seq}
\end{figure}

% ------------------------------------------------------------
\section{Design decisions and alternatives considered}
\label{sec:design-decisions}

Table~\ref{fig:design-decisions} summarises the ten decisions that
realise the specification, names the alternative considered for
each, and records the specific technical reason for rejection. The
rest of this chapter expands on the decisions that drive the
experiment outcome; the remaining rows are covered by the table.

\begin{figure}[ht]
    \centering
    \includegraphics[width=\linewidth]{Chapters/figures/fig-5-4-decisions-table.png}
    \caption{Design decisions, rejected alternatives and
    traceability to FR, NFR and ER classes.}
    \label{fig:design-decisions}
\end{figure}

\paragraph{Modality-isolation boundary.} The single branch point
keeps ER1 mechanical rather than editorial. Two parallel pipelines
would satisfy the letter of ER1 on day one and violate it silently
on any later change to, say, the prompt, the post-processor or the
JSON parser. The chosen shape makes provenance of any difference
directly inspectable: if the shared input and shared output are the
same byte-for-byte, the difference must come from the branch.

\paragraph{Cross-provider symmetry.} ER7 requires that swapping a
provider cannot alter pipeline code. The \texttt{LLMProvider} ABC
plus registry in \S\ref{sec:design-components} satisfies this
structurally: the pipeline sees only the ABC, so every provider is
exchanged at the level of a config key. The alternative was a
single adapter class with \texttt{if provider == "openai"}
branches; it was rejected because every provider-specific test
would then need to touch pipeline logic, and a future third
provider would require editing the pipeline rather than adding a
file.

\paragraph{Null-aware classification.} ER6 requires the four
classes of the taxonomy defined in \S\ref{sec:spec-evaluation}. The
classification does not live only in the evaluator: it shapes the
failure-diagnosis heuristic in \S\ref{sec:design-evalframework} and
the reporting in Chapter~\ref{chap:evaluation}. The rejected
alternative was a binary exact-match with both-null counted as
correct; it was rejected because on sparse bills (where many
fields are legitimately absent) it can return above 80\% accuracy
on a model that outputs \texttt{null} for every field, which is the
opposite of the discrimination the RQs require.

\paragraph{JSON parse hardening.} Providers do not agree on output
format (\S\ref{sec:lr-eval}). OpenAI Structured Outputs guarantees
a schema-valid JSON object. Anthropic without Structured Outputs
returns text that may be wrapped in markdown fences and may contain
preamble. A retry-with-repair prompt would satisfy this trivially
by asking the model to reformat, but at the cost of changing the
prompt between conditions and breaking ER2. The chosen design is
a four-stage offline fallback (\S\ref{sec:design-execution}) that
never reissues a call.

% ------------------------------------------------------------
\section{Pipeline architecture}
\label{sec:design-pipeline}

\subsection{Shared preprocessing}
\label{sec:design-pipeline-preproc}

Every document enters a single rasteriser before the branch, so the
pixel source for both paths is bit-identical. The alternative would
be to rasterise once per path; it was rejected because PDF
rasterisation is non-deterministic across invocations in the
general case (font hinting, anti-aliasing) and two independent
rasterisations would introduce a second unintended variable.

\subsection{Modality branch}

The OCR path runs Tesseract with the language taken from the
manifest; only the recognised text is forwarded to the LLM. The
vision path sends the first two rasterised pages directly to the
LLM in vision mode. The two-page cap follows from the scope
commitment in \S\ref{sec:spec-scope}: most target fields appear on
the first page or page two. The rejected alternative was to pass
every page; it was rejected because vision tokens scale
near-linearly with page area and the cost would grow quadratically
with document size for no demonstrated accuracy benefit on target
fields.

\subsection{Shared post-processing}

Every LLM response passes through the same four-stage chain: JSON
parse, field-level normalisation, Pydantic validation, per-document
persistence. Parsing and validation are separated so that a parse
failure is distinguishable from a schema violation: the former
keeps the raw output and retries offline
(\S\ref{sec:design-execution}), the latter fails loudly at
evaluation time. Pydantic is chosen over raw \texttt{jsonschema}
because the same model serves as the OpenAI Structured Outputs
schema, the LLM-output container, the normalisation target, the
evaluator input and the reporting envelope; a single Python class
removes five format-translation layers that would otherwise need to
be kept in sync.

% ------------------------------------------------------------
\section{Component design}
\label{sec:design-components}

Figure~\ref{fig:design-strategy} shows the two strategy hierarchies
and the registry that binds them. Every component visible to the
pipeline is an ABC; every concrete implementation is opaque to the
pipeline and instantiated at run time by
\texttt{registry.get\_provider(config)} (lazy import).

\begin{figure}[ht]
    \centering
    \includegraphics[width=\linewidth]{Chapters/figures/fig-5-3-strategy-pattern.png}
    \caption{Provider and OCR strategy hierarchies bound at
    run time by the registry. The rejected alternative
    (hard-coded dispatch) is shown beneath.}
    \label{fig:design-strategy}
\end{figure}

\subsection{LLM provider abstraction}
\label{sec:design-components-llm}

\texttt{LLMProvider} exposes \texttt{extract\_from\_text},
\texttt{extract\_from\_image} and \texttt{get\_model\_id}. Both
extraction methods return a uniform triple
(\texttt{raw\_output}, \texttt{token\_usage}, \texttt{latency\_ms}).
The pipeline calls only the ABC; provider identity enters through
the registry and leaves in the persisted \texttt{model\_id} so that
downstream evaluation can slice by model. The rejected alternative
was a hard-coded \texttt{if}/\texttt{elif} dispatch inside
\texttt{pipeline.py}; its rejection reason is recorded in
Figure~\ref{fig:design-decisions}.

\subsection{OCR engine abstraction}
\label{sec:design-components-ocr}

\texttt{OCREngine} exposes \texttt{extract\_text(image, lang)} and
is implemented by \texttt{TesseractOCR}. Tesseract is the chosen
concrete engine because it executes locally (no confound with a
second paid API in the OCR path) and ships language packs for the
four languages in scope. The abstraction survives the choice: a
later experiment can substitute EasyOCR or PaddleOCR by adding a
module without touching the pipeline, which is what ER5's
failure-attribution story needs to remain falsifiable.

\subsection{Prompt system}
\label{sec:design-components-prompt}

Prompts are versioned files in \texttt{prompts/}, selected by
\texttt{llm.prompt\_file}. The schema description is not written
into the prompt; it is generated from
\texttt{BillExtraction.schema\_description()} and injected through
the \texttt{\{schema\_description\}} placeholder at render time.
The rendered prompt is archived per run (\S\ref{sec:design-execution}) so a diff-only audit proves prompt
identity across conditions, which is what ER2 requires. Inline
string constants were rejected because differences between
conditions cannot be audited by reading a file.

\subsection{Provider-API asymmetry}
\label{sec:design-components-api}

The two providers use different APIs and different mechanisms for
structured output, and the design records the asymmetry rather
than papering over it. OpenAI is called through the Responses API
with Structured Outputs: \texttt{text\_format=BillExtraction}
constrains decoding to schema-valid JSON. Anthropic is called
through the Messages API without Structured Outputs; the same
prompt instructs a JSON response, and the provider's tendency to
wrap JSON in markdown fences is handled by
\texttt{\_strip\_json\_fencing()} before the shared parser runs.

This asymmetry is deliberate. Forcing both providers into
prompt-instructed JSON for the sake of parity would discard the
capability that gives OpenAI guaranteed schema adherence and shift
the risk profile of the experiment silently. The chosen design
preserves each provider's strongest available output mode and
records the raw LLM response per document (NFR8) so that format
failure is separable from content failure in post-hoc analysis. The
evaluation chapter revisits this asymmetry explicitly; a design
choice that hides it would make that analysis impossible.

\subsection{Configuration and CLI}

The runtime is configured by a YAML file (\texttt{config/default.yaml}) with CLI deep-merge overrides
for the handful of flags exercised per run (\texttt{--mode},
\texttt{--provider}, \texttt{--model}, \texttt{--force},
\texttt{--manifest}). The effective configuration after merging is
archived with the run. An environment-variable-only alternative was
rejected because at run-time audit it lacks the single-file
self-description that reproducibility (ER4) assumes.

% ------------------------------------------------------------
\section{Data model}
\label{sec:design-datamodel}

One Pydantic class, \texttt{BillExtraction}
(Table~\ref{fig:schema}), is the contract between prompt, LLM
output, normaliser, evaluator and reporter. Every field is
\texttt{Optional}, which is not a defensive default but the shape
the 4-class null taxonomy assumes: a distinct \texttt{null} is the
only way \texttt{hallucination} and \texttt{omission} can be
separated from \texttt{incorrect}. \texttt{PipelineResult} bundles
the extraction with the artefacts the evaluator and the
cost/latency accounting need (raw LLM output, OCR text, timings,
token usage, model id, pipeline mode); it is not a database row.
\texttt{DocumentEntry} is the resolved manifest row with its file
paths already verified; only entries that pass the six-step manifest
validation (\S\ref{sec:spec-experimental}) reach the pipeline.

Per-layer dataclasses were rejected because every field rename
would then be a five-file change (prompt description, LLM output
parser, normaliser dispatch, evaluator comparison, report column
header) with silent divergence if one file is missed; a single
Pydantic class fails the rename loudly.

% ------------------------------------------------------------
\section{Normalisation}
\label{sec:design-normalisation}

The same normaliser runs over predictions and ground truth before
any comparison. This is the one design choice that most directly
protects the meaning of ``correct'' against surface variation:
\texttt{"\texteuro"} and \texttt{"EUR"}, \texttt{"01/03/2024"} and
\texttt{"2024-03-01"}, \texttt{"OVO Energy Ltd"} and \texttt{"OVO
Energy"} all collapse to the same value before they are compared.
Normalising predictions alone was rejected because it would push
format parity onto the annotators: the set of correct surface forms
would be whatever the normaliser happened to recognise, not the set
of meaningfully correct forms. That trade breaks ER6.

Normalisation dispatches through \texttt{FIELD\_NORMALISERS}, a
mapping from field name to normaliser function. A monolithic
\texttt{normalise(fields)} with a chain of \texttt{if} statements
was rejected because adding a new field would then require editing
the normaliser, the comparator and the reporter in three places;
the dispatch table localises the addition to one entry in one table.
The float parser is the one non-trivial normaliser: it distinguishes
EU and US decimal conventions from the positions of the group and
decimal separators, falling back to standard parsing when both are
absent. LLM-as-judge comparison was rejected as an alternative
because the evaluator becomes non-reproducible the moment it calls
a remote model, breaking ER4.

% ------------------------------------------------------------
\section{Evaluation framework}
\label{sec:design-evalframework}

The evaluation is specified declaratively in
\S\ref{sec:spec-evaluation}; this section states its decision-tree
realisation.

\subsection{Metric set}

Field-level accuracy (over documents), document-level accuracy
(fraction of documents with no incorrect, hallucinated or omitted
field), overall accuracy (all fields across all documents), and
slice accuracy (by language, by utility type) are the four reported
metrics. The split addresses the failure mode named in
\S\ref{sec:lr-eval}: a single aggregate number can hide systematic
language-specific degradation, and coarse-grained metrics that
aggregate across a 12-field schema can hide field-class-specific
degradation (for example, a pipeline that is accurate on
\texttt{provider\_name} but randomly wrong on \texttt{due\_date}).
The six canonical charts listed in \S\ref{sec:spec-evaluation} are
each pinned to a specific research question. Overall-accuracy and
per-field-accuracy answer RQ1 (modality). Per-stage timing and
cost-per-document answer RQ2 (cost and latency).
Accuracy-by-language and accuracy-by-utility type answer RQ3
(multilingual and cross-issuer degradation). No chart exists
without a question.

\subsection{Null-aware classification}

Each field-level comparison emits one of five classifications
(Table~\ref{fig:null-taxonomy}): \texttt{correct},
\texttt{incorrect}, \texttt{hallucination}, \texttt{omission},
\texttt{both-null}. The first four are counted in the accuracy
denominator; \texttt{both-null} is counted as a match but excluded
from the denominator. The exclusion is what prevents a
\texttt{null}-emitting pipeline from scoring highly on sparse
bills. The classification is decided at compare time and persisted
so that the failure-diagnosis heuristic below can operate on it
without re-running comparison logic.

\subsection{Failure diagnosis}
\label{sec:design-evalframework-diagnosis}

For each non-correct field in an OCR-mediated condition, the
heuristic reads the stored \texttt{ocr\_text.txt} and searches it,
case-insensitively, for the ground-truth value. A hit means
Tesseract captured the value and the LLM misused it: the failure
is attributed to the LLM reasoning step. A miss means Tesseract did
not capture the value and no downstream stage could have succeeded:
the failure is attributed to OCR. In vision-mediated conditions no
intermediate text artefact exists, so all non-correct fields are
attributed to the vision model; the asymmetry is recorded and
discussed in Chapter~\ref{chap:evaluation} rather than hidden behind
a weaker heuristic.

Training a classifier for failure attribution was rejected because
it makes the evaluator itself an unvalidated ML component:
disagreements between the classifier and ground truth would then
need their own taxonomy, and an error in the attribution classifier
would be indistinguishable from an error in the pipeline under
test. The substring heuristic is inspectable; its known
false-positive modes (repeated digits on dates) are recorded in
Chapter~\ref{chap:evaluation}.

\subsection{Cross-run comparator}

\texttt{compare\_runs()} reads per-run \texttt{evaluation.json}
files and assembles the accuracy matrix, performance matrix, null
analysis, and slice comparisons for every run under comparison.
The output feeds the factorial analysis specified in
\S\ref{sec:spec-experimental}: the four cells of the $2 \times 2$
matrix are four runs to the comparator, and the comparator is the
only place accuracy across runs is aggregated. No comparison is
ever computed inside the pipeline, so pipeline changes cannot
accidentally rewrite comparison logic.

% ------------------------------------------------------------
\section{Execution model and reproducibility}
\label{sec:design-execution}

Figure~\ref{fig:design-rundir} gives the per-run directory layout.
A run produces one directory, named
\texttt{\{timestamp\}\_\{provider\}\_\{model\}\_\{mode\}}, that
contains the effective configuration, the manifest snapshot at run
time, the exact prompt template, a summary, and a
\texttt{documents/} subdirectory with one folder per document.
Every document folder holds the normalised extraction, the raw LLM
response, the OCR text when applicable, per-stage timings, token
and cost metadata, and on failure an \texttt{error.json} with the
full traceback. The evaluator is a consumer of this directory: it
reads only from disk, never from in-memory objects owned by the
pipeline, which is why pipeline and evaluator are independently
runnable.

\begin{figure}[ht]
    \centering
    \includegraphics[width=0.9\linewidth]{Chapters/figures/fig-5-6-run-directory.png}
    \caption{Per-run directory layout. The evaluator reads only
    from this directory; pipeline and evaluator never share memory,
    so each can be run, replayed and inspected independently.}
    \label{fig:design-rundir}
\end{figure}

Error isolation is mechanical rather than editorial: each document
is processed inside a \texttt{try}/\texttt{except}, and on failure
an \texttt{error.json} with full traceback is written and the run
proceeds. The evaluator skips any document with an
\texttt{error.json}. A crash on document forty therefore loses one
document rather than forty-one.

Figure~\ref{fig:design-parse} gives the JSON parse hardening chain.
The chain is local and offline: each stage rewrites the input
locally and retries the parse, and the final stage raises with a
diagnostic excerpt of the raw output (already archived for the
run). The rejected alternative was a retry-with-repair round-trip
to the model; it was rejected because it mutates the prompt
conditional on provider behaviour, and any prompt mutation breaks
ER2.

\begin{figure}[ht]
    \centering
    \includegraphics[width=0.75\linewidth]{Chapters/figures/fig-5-5-json-parse-fallback.png}
    \caption{JSON parse hardening chain. Every stage is local and
    offline; no stage reissues a call to the model.}
    \label{fig:design-parse}
\end{figure}

Resume is the default: if a document already has an
\texttt{extraction.json}, it is skipped. The \texttt{--force} flag
re-runs everything. This makes partial re-runs cheap, which matters
because API cost per condition is non-trivial (\S\ref{sec:spec-experimental}).

% ------------------------------------------------------------
\section{Maintenance considerations}
\label{sec:design-maintenance}

Five dependency risks are named rather than implied. Poppler
(binary dependency of \texttt{pdf2image}) and Tesseract (with the
four language packs in scope) are system-level dependencies that
must be present on every run host; the \texttt{README} and
\texttt{requirements.txt} list both. Provider SDKs
(\texttt{openai}, \texttt{anthropic}) are pinned in
\texttt{requirements.txt} because both providers make incompatible
API changes between major versions. Pydantic is pinned to v2
because OpenAI Structured Outputs uses v2-specific schema
metadata. The chart backend (\texttt{matplotlib}) is pinned loosely
because its failure mode is visual, not semantic.

Three extensibility points cover the expected directions of
future change. Adding a provider requires one new module and one
entry in \texttt{registry.\_PROVIDERS}. Adding an OCR engine
requires one new \texttt{OCREngine} subclass. Adding a schema field
requires one Pydantic attribute plus one entry in
\texttt{FIELD\_NORMALISERS}; the evaluator and reporter iterate
over \texttt{BillExtraction.SCHEMA\_FIELDS}, so they pick up new
fields automatically.

\paragraph{Expected deviations.} Three deviations from this design
are anticipated and will be recorded in Chapter~\ref{chap:impl}~\S\ref{sec:impl-deviations}.
First, the provider-API asymmetry in
\S\ref{sec:design-components-api} is itself a deviation from the
na{\"i}ve design of ``one call shape per provider''; the deviation
is documented here because hiding it would invalidate the
evaluation. Second, some providers may not return
\texttt{latency\_ms} uniformly across SDK versions, and the Timer
wrap in \texttt{src/performance.py} is expected to supersede the
provider field. Third, Tesseract language packs are not
universally installed on arbitrary hosts, and a missing pack
downgrades an OCR run to empty text; this is a ground-truth issue
for that condition, not a bug, and will be handled by the
failure-attribution heuristic.

% ------------------------------------------------------------
\section{Traceability summary}
\label{sec:design-traceability}

Figure~\ref{fig:design-trace} gives the end-to-end chain for each
experimental-requirement class. Every row connects a literature
gap to a requirement, to the design realisation in this chapter,
to the test or evaluation hook that will exercise it. The table is
the single defensible artefact for Trait~5: a marker can verify
the chain without reading the chapters in sequence.

\begin{figure}[ht]
    \centering
    \includegraphics[width=\linewidth]{Chapters/figures/fig-5-7-traceability-matrix.png}
    \caption{End-to-end traceability from literature gap to
    evaluation hook for every experimental requirement.}
    \label{fig:design-trace}
\end{figure}

% ------------------------------------------------------------
\section{Summary and forward bridge}
\label{sec:design-summary}

This chapter turned each requirement of Chapter~\ref{chap:spec}
into a design decision, named an alternative for each decision, and
stated the technical reason for rejection. The single architectural
commitment that dominates the chapter is the modality-isolation
boundary: every other decision exists to make that boundary
auditable by someone other than the author.
Chapter~\ref{chap:impl} takes this design and builds it, recording
the real deviations encountered against the expected deviations in
\S\ref{sec:design-maintenance}. Chapter~\ref{chap:evaluation} tests
whether the design produces the results predicted by the
hypotheses in \S\ref{sec:spec-objectives}, including the
alternative-design comparison that returns to the rejected rows of
Figure~\ref{fig:design-decisions} and evaluates them against the
same data.
.
%  Do NOT add \documentclass or \begin{document} here.
%  Follows Chapter 4 (Specification); precedes Chapter 6 (Implementation).
%  Scope: architecture, design decisions with named alternatives and
%  rejection reasons, pipeline and component design, data model,
%  normalisation, evaluation framework, execution model, maintenance,
%  traceability. Written before implementation; "expected deviations"
%  section flags constraints that will shape Chapter 6.
% ============================================================

\chapter{Design}
\label{chap:design}

% ------------------------------------------------------------
\section{From specification to design}
\label{sec:design-bridge}

Chapter~\ref{chap:spec} fixed what the system must satisfy;
this chapter states how it is realised and what was rejected on the
way. Every design decision below is anchored in a requirement from
\S\ref{sec:spec-fr}--\S\ref{sec:spec-experimental-requirements} and
traces upstream to a gap named in \S\ref{sec:lr-gap}. Four
requirements shape the chapter. ER1 (modality isolation) dictates
that the two pipelines share every component except the one under
test, which constrains the architectural boundary in \S\ref{sec:design-overview}. ER2 (multilingual and prompt-controlled
coverage) dictates that the prompt and the LLM are auditable
artefacts, not inlined strings, which forces the component
abstractions in \S\ref{sec:design-components}. ER6 (null-aware
evaluation) dictates that both predictions and ground truth are
normalised identically before comparison, which drives the
normalisation design in \S\ref{sec:design-normalisation}. ER7
(cross-provider generality) dictates that swapping a provider cannot
change pipeline code, which drives the registry in
\S\ref{sec:design-components}. The remaining decisions fall out from
these four.

% ------------------------------------------------------------
\section{System overview}
\label{sec:design-overview}

Figure~\ref{fig:design-arch} gives the architecture. Three
boundaries matter. \textit{Shared input}
(manifest, loader, PDF, rasteriser) is identical across every
condition, so any difference in raw pixel content is a data-source
effect rather than a pipeline effect. \textit{Modality branch}
contains the single step that varies: either Tesseract followed by
the LLM in text mode, or the LLM in vision mode. \textit{Shared
output} (JSON parse, normalise, validate, persist, evaluate) is
identical across every condition, so any difference in final output
distribution is attributable to the branch content rather than to
the finalisation logic.

The branch is the only place ER1 can be violated, which is why it
is drawn as a single box rather than as two parallel pipelines: two
parallel pipelines have two places to drift. Holding the boundary
tight enough for one provider to be swapped independently of one
modality (\S\ref{sec:design-components}) costs some visual symmetry
in the figure but removes the dominant confound identified in
\S\ref{sec:lr-controlled}.

\begin{figure}[ht]
    \centering
    \includegraphics[width=\linewidth]{Chapters/figures/fig-5-1-architecture.png}
    \caption{End-to-end architecture. Shaded regions are held
    constant across conditions; the central dashed region is the
    modality branch where the input differs. The design isolates
    the ER1 control to a single step rather than to a whole
    pipeline.}
    \label{fig:design-arch}
\end{figure}

Figure~\ref{fig:design-seq} refines the same picture at the
per-document scale. The two paths differ only in whether Tesseract
sits between the rasteriser and the LLM; every call on either side
of that step is identical. The timing keys in brackets match the
per-stage measurements written to \texttt{timings.json} and drive
the latency analysis for RQ2.

\begin{figure}[ht]
    \centering
    \includegraphics[width=\linewidth]{Chapters/figures/fig-5-2-sequence-diagram.png}
    \caption{Per-document sequence for both modalities. The only
    step that varies is whether Tesseract sits in the chain.}
    \label{fig:design-seq}
\end{figure}

% ------------------------------------------------------------
\section{Design decisions and alternatives considered}
\label{sec:design-decisions}

Table~\ref{fig:design-decisions} summarises the ten decisions that
realise the specification, names the alternative considered for
each, and records the specific technical reason for rejection. The
rest of this chapter expands on the decisions that drive the
experiment outcome; the remaining rows are covered by the table.

\begin{figure}[ht]
    \centering
    \includegraphics[width=\linewidth]{Chapters/figures/fig-5-4-decisions-table.png}
    \caption{Design decisions, rejected alternatives and
    traceability to FR, NFR and ER classes.}
    \label{fig:design-decisions}
\end{figure}

\paragraph{Modality-isolation boundary.} The single branch point
keeps ER1 mechanical rather than editorial. Two parallel pipelines
would satisfy the letter of ER1 on day one and violate it silently
on any later change to, say, the prompt, the post-processor or the
JSON parser. The chosen shape makes provenance of any difference
directly inspectable: if the shared input and shared output are the
same byte-for-byte, the difference must come from the branch.

\paragraph{Cross-provider symmetry.} ER7 requires that swapping a
provider cannot alter pipeline code. The \texttt{LLMProvider} ABC
plus registry in \S\ref{sec:design-components} satisfies this
structurally: the pipeline sees only the ABC, so every provider is
exchanged at the level of a config key. The alternative was a
single adapter class with \texttt{if provider == "openai"}
branches; it was rejected because every provider-specific test
would then need to touch pipeline logic, and a future third
provider would require editing the pipeline rather than adding a
file.

\paragraph{Null-aware classification.} ER6 requires the four
classes of the taxonomy defined in \S\ref{sec:spec-evaluation}. The
classification does not live only in the evaluator: it shapes the
failure-diagnosis heuristic in \S\ref{sec:design-evalframework} and
the reporting in Chapter~\ref{chap:evaluation}. The rejected
alternative was a binary exact-match with both-null counted as
correct; it was rejected because on sparse bills (where many
fields are legitimately absent) it can return above 80\% accuracy
on a model that outputs \texttt{null} for every field, which is the
opposite of the discrimination the RQs require.

\paragraph{JSON parse hardening.} Providers do not agree on output
format (\S\ref{sec:lr-eval}). OpenAI Structured Outputs guarantees
a schema-valid JSON object. Anthropic without Structured Outputs
returns text that may be wrapped in markdown fences and may contain
preamble. A retry-with-repair prompt would satisfy this trivially
by asking the model to reformat, but at the cost of changing the
prompt between conditions and breaking ER2. The chosen design is
a four-stage offline fallback (\S\ref{sec:design-execution}) that
never reissues a call.

% ------------------------------------------------------------
\section{Pipeline architecture}
\label{sec:design-pipeline}

\subsection{Shared preprocessing}
\label{sec:design-pipeline-preproc}

Every document enters a single rasteriser before the branch, so the
pixel source for both paths is bit-identical. The alternative would
be to rasterise once per path; it was rejected because PDF
rasterisation is non-deterministic across invocations in the
general case (font hinting, anti-aliasing) and two independent
rasterisations would introduce a second unintended variable.

\subsection{Modality branch}

The OCR path runs Tesseract with the language taken from the
manifest; only the recognised text is forwarded to the LLM. The
vision path sends the first two rasterised pages directly to the
LLM in vision mode. The two-page cap follows from the scope
commitment in \S\ref{sec:spec-scope}: most target fields appear on
the first page or page two. The rejected alternative was to pass
every page; it was rejected because vision tokens scale
near-linearly with page area and the cost would grow quadratically
with document size for no demonstrated accuracy benefit on target
fields.

\subsection{Shared post-processing}

Every LLM response passes through the same four-stage chain: JSON
parse, field-level normalisation, Pydantic validation, per-document
persistence. Parsing and validation are separated so that a parse
failure is distinguishable from a schema violation: the former
keeps the raw output and retries offline
(\S\ref{sec:design-execution}), the latter fails loudly at
evaluation time. Pydantic is chosen over raw \texttt{jsonschema}
because the same model serves as the OpenAI Structured Outputs
schema, the LLM-output container, the normalisation target, the
evaluator input and the reporting envelope; a single Python class
removes five format-translation layers that would otherwise need to
be kept in sync.

% ------------------------------------------------------------
\section{Component design}
\label{sec:design-components}

Figure~\ref{fig:design-strategy} shows the two strategy hierarchies
and the registry that binds them. Every component visible to the
pipeline is an ABC; every concrete implementation is opaque to the
pipeline and instantiated at run time by
\texttt{registry.get\_provider(config)} (lazy import).

\begin{figure}[ht]
    \centering
    \includegraphics[width=\linewidth]{Chapters/figures/fig-5-3-strategy-pattern.png}
    \caption{Provider and OCR strategy hierarchies bound at
    run time by the registry. The rejected alternative
    (hard-coded dispatch) is shown beneath.}
    \label{fig:design-strategy}
\end{figure}

\subsection{LLM provider abstraction}
\label{sec:design-components-llm}

\texttt{LLMProvider} exposes \texttt{extract\_from\_text},
\texttt{extract\_from\_image} and \texttt{get\_model\_id}. Both
extraction methods return a uniform triple
(\texttt{raw\_output}, \texttt{token\_usage}, \texttt{latency\_ms}).
The pipeline calls only the ABC; provider identity enters through
the registry and leaves in the persisted \texttt{model\_id} so that
downstream evaluation can slice by model. The rejected alternative
was a hard-coded \texttt{if}/\texttt{elif} dispatch inside
\texttt{pipeline.py}; its rejection reason is recorded in
Figure~\ref{fig:design-decisions}.

\subsection{OCR engine abstraction}
\label{sec:design-components-ocr}

\texttt{OCREngine} exposes \texttt{extract\_text(image, lang)} and
is implemented by \texttt{TesseractOCR}. Tesseract is the chosen
concrete engine because it executes locally (no confound with a
second paid API in the OCR path) and ships language packs for the
four languages in scope. The abstraction survives the choice: a
later experiment can substitute EasyOCR or PaddleOCR by adding a
module without touching the pipeline, which is what ER5's
failure-attribution story needs to remain falsifiable.

\subsection{Prompt system}
\label{sec:design-components-prompt}

Prompts are versioned files in \texttt{prompts/}, selected by
\texttt{llm.prompt\_file}. The schema description is not written
into the prompt; it is generated from
\texttt{BillExtraction.schema\_description()} and injected through
the \texttt{\{schema\_description\}} placeholder at render time.
The rendered prompt is archived per run (\S\ref{sec:design-execution}) so a diff-only audit proves prompt
identity across conditions, which is what ER2 requires. Inline
string constants were rejected because differences between
conditions cannot be audited by reading a file.

\subsection{Provider-API asymmetry}
\label{sec:design-components-api}

The two providers use different APIs and different mechanisms for
structured output, and the design records the asymmetry rather
than papering over it. OpenAI is called through the Responses API
with Structured Outputs: \texttt{text\_format=BillExtraction}
constrains decoding to schema-valid JSON. Anthropic is called
through the Messages API without Structured Outputs; the same
prompt instructs a JSON response, and the provider's tendency to
wrap JSON in markdown fences is handled by
\texttt{\_strip\_json\_fencing()} before the shared parser runs.

This asymmetry is deliberate. Forcing both providers into
prompt-instructed JSON for the sake of parity would discard the
capability that gives OpenAI guaranteed schema adherence and shift
the risk profile of the experiment silently. The chosen design
preserves each provider's strongest available output mode and
records the raw LLM response per document (NFR8) so that format
failure is separable from content failure in post-hoc analysis. The
evaluation chapter revisits this asymmetry explicitly; a design
choice that hides it would make that analysis impossible.

\subsection{Configuration and CLI}

The runtime is configured by a YAML file (\texttt{config/default.yaml}) with CLI deep-merge overrides
for the handful of flags exercised per run (\texttt{--mode},
\texttt{--provider}, \texttt{--model}, \texttt{--force},
\texttt{--manifest}). The effective configuration after merging is
archived with the run. An environment-variable-only alternative was
rejected because at run-time audit it lacks the single-file
self-description that reproducibility (ER4) assumes.

% ------------------------------------------------------------
\section{Data model}
\label{sec:design-datamodel}

One Pydantic class, \texttt{BillExtraction}
(Table~\ref{fig:schema}), is the contract between prompt, LLM
output, normaliser, evaluator and reporter. Every field is
\texttt{Optional}, which is not a defensive default but the shape
the 4-class null taxonomy assumes: a distinct \texttt{null} is the
only way \texttt{hallucination} and \texttt{omission} can be
separated from \texttt{incorrect}. \texttt{PipelineResult} bundles
the extraction with the artefacts the evaluator and the
cost/latency accounting need (raw LLM output, OCR text, timings,
token usage, model id, pipeline mode); it is not a database row.
\texttt{DocumentEntry} is the resolved manifest row with its file
paths already verified; only entries that pass the six-step manifest
validation (\S\ref{sec:spec-experimental}) reach the pipeline.

Per-layer dataclasses were rejected because every field rename
would then be a five-file change (prompt description, LLM output
parser, normaliser dispatch, evaluator comparison, report column
header) with silent divergence if one file is missed; a single
Pydantic class fails the rename loudly.

% ------------------------------------------------------------
\section{Normalisation}
\label{sec:design-normalisation}

The same normaliser runs over predictions and ground truth before
any comparison. This is the one design choice that most directly
protects the meaning of ``correct'' against surface variation:
\texttt{"\texteuro"} and \texttt{"EUR"}, \texttt{"01/03/2024"} and
\texttt{"2024-03-01"}, \texttt{"OVO Energy Ltd"} and \texttt{"OVO
Energy"} all collapse to the same value before they are compared.
Normalising predictions alone was rejected because it would push
format parity onto the annotators: the set of correct surface forms
would be whatever the normaliser happened to recognise, not the set
of meaningfully correct forms. That trade breaks ER6.

Normalisation dispatches through \texttt{FIELD\_NORMALISERS}, a
mapping from field name to normaliser function. A monolithic
\texttt{normalise(fields)} with a chain of \texttt{if} statements
was rejected because adding a new field would then require editing
the normaliser, the comparator and the reporter in three places;
the dispatch table localises the addition to one entry in one table.
The float parser is the one non-trivial normaliser: it distinguishes
EU and US decimal conventions from the positions of the group and
decimal separators, falling back to standard parsing when both are
absent. LLM-as-judge comparison was rejected as an alternative
because the evaluator becomes non-reproducible the moment it calls
a remote model, breaking ER4.

% ------------------------------------------------------------
\section{Evaluation framework}
\label{sec:design-evalframework}

The evaluation is specified declaratively in
\S\ref{sec:spec-evaluation}; this section states its decision-tree
realisation.

\subsection{Metric set}

Field-level accuracy (over documents), document-level accuracy
(fraction of documents with no incorrect, hallucinated or omitted
field), overall accuracy (all fields across all documents), and
slice accuracy (by language, by utility type) are the four reported
metrics. The split addresses the failure mode named in
\S\ref{sec:lr-eval}: a single aggregate number can hide systematic
language-specific degradation, and coarse-grained metrics that
aggregate across a 12-field schema can hide field-class-specific
degradation (for example, a pipeline that is accurate on
\texttt{provider\_name} but randomly wrong on \texttt{due\_date}).
The six canonical charts listed in \S\ref{sec:spec-evaluation} are
each pinned to a specific research question. Overall-accuracy and
per-field-accuracy answer RQ1 (modality). Per-stage timing and
cost-per-document answer RQ2 (cost and latency).
Accuracy-by-language and accuracy-by-utility type answer RQ3
(multilingual and cross-issuer degradation). No chart exists
without a question.

\subsection{Null-aware classification}

Each field-level comparison emits one of five classifications
(Table~\ref{fig:null-taxonomy}): \texttt{correct},
\texttt{incorrect}, \texttt{hallucination}, \texttt{omission},
\texttt{both-null}. The first four are counted in the accuracy
denominator; \texttt{both-null} is counted as a match but excluded
from the denominator. The exclusion is what prevents a
\texttt{null}-emitting pipeline from scoring highly on sparse
bills. The classification is decided at compare time and persisted
so that the failure-diagnosis heuristic below can operate on it
without re-running comparison logic.

\subsection{Failure diagnosis}
\label{sec:design-evalframework-diagnosis}

For each non-correct field in an OCR-mediated condition, the
heuristic reads the stored \texttt{ocr\_text.txt} and searches it,
case-insensitively, for the ground-truth value. A hit means
Tesseract captured the value and the LLM misused it: the failure
is attributed to the LLM reasoning step. A miss means Tesseract did
not capture the value and no downstream stage could have succeeded:
the failure is attributed to OCR. In vision-mediated conditions no
intermediate text artefact exists, so all non-correct fields are
attributed to the vision model; the asymmetry is recorded and
discussed in Chapter~\ref{chap:evaluation} rather than hidden behind
a weaker heuristic.

Training a classifier for failure attribution was rejected because
it makes the evaluator itself an unvalidated ML component:
disagreements between the classifier and ground truth would then
need their own taxonomy, and an error in the attribution classifier
would be indistinguishable from an error in the pipeline under
test. The substring heuristic is inspectable; its known
false-positive modes (repeated digits on dates) are recorded in
Chapter~\ref{chap:evaluation}.

\subsection{Cross-run comparator}

\texttt{compare\_runs()} reads per-run \texttt{evaluation.json}
files and assembles the accuracy matrix, performance matrix, null
analysis, and slice comparisons for every run under comparison.
The output feeds the factorial analysis specified in
\S\ref{sec:spec-experimental}: the four cells of the $2 \times 2$
matrix are four runs to the comparator, and the comparator is the
only place accuracy across runs is aggregated. No comparison is
ever computed inside the pipeline, so pipeline changes cannot
accidentally rewrite comparison logic.

% ------------------------------------------------------------
\section{Execution model and reproducibility}
\label{sec:design-execution}

Figure~\ref{fig:design-rundir} gives the per-run directory layout.
A run produces one directory, named
\texttt{\{timestamp\}\_\{provider\}\_\{model\}\_\{mode\}}, that
contains the effective configuration, the manifest snapshot at run
time, the exact prompt template, a summary, and a
\texttt{documents/} subdirectory with one folder per document.
Every document folder holds the normalised extraction, the raw LLM
response, the OCR text when applicable, per-stage timings, token
and cost metadata, and on failure an \texttt{error.json} with the
full traceback. The evaluator is a consumer of this directory: it
reads only from disk, never from in-memory objects owned by the
pipeline, which is why pipeline and evaluator are independently
runnable.

\begin{figure}[ht]
    \centering
    \includegraphics[width=0.9\linewidth]{Chapters/figures/fig-5-6-run-directory.png}
    \caption{Per-run directory layout. The evaluator reads only
    from this directory; pipeline and evaluator never share memory,
    so each can be run, replayed and inspected independently.}
    \label{fig:design-rundir}
\end{figure}

Error isolation is mechanical rather than editorial: each document
is processed inside a \texttt{try}/\texttt{except}, and on failure
an \texttt{error.json} with full traceback is written and the run
proceeds. The evaluator skips any document with an
\texttt{error.json}. A crash on document forty therefore loses one
document rather than forty-one.

Figure~\ref{fig:design-parse} gives the JSON parse hardening chain.
The chain is local and offline: each stage rewrites the input
locally and retries the parse, and the final stage raises with a
diagnostic excerpt of the raw output (already archived for the
run). The rejected alternative was a retry-with-repair round-trip
to the model; it was rejected because it mutates the prompt
conditional on provider behaviour, and any prompt mutation breaks
ER2.

\begin{figure}[ht]
    \centering
    \includegraphics[width=0.75\linewidth]{Chapters/figures/fig-5-5-json-parse-fallback.png}
    \caption{JSON parse hardening chain. Every stage is local and
    offline; no stage reissues a call to the model.}
    \label{fig:design-parse}
\end{figure}

Resume is the default: if a document already has an
\texttt{extraction.json}, it is skipped. The \texttt{--force} flag
re-runs everything. This makes partial re-runs cheap, which matters
because API cost per condition is non-trivial (\S\ref{sec:spec-experimental}).

% ------------------------------------------------------------
\section{Maintenance considerations}
\label{sec:design-maintenance}

Five dependency risks are named rather than implied. Poppler
(binary dependency of \texttt{pdf2image}) and Tesseract (with the
four language packs in scope) are system-level dependencies that
must be present on every run host; the \texttt{README} and
\texttt{requirements.txt} list both. Provider SDKs
(\texttt{openai}, \texttt{anthropic}) are pinned in
\texttt{requirements.txt} because both providers make incompatible
API changes between major versions. Pydantic is pinned to v2
because OpenAI Structured Outputs uses v2-specific schema
metadata. The chart backend (\texttt{matplotlib}) is pinned loosely
because its failure mode is visual, not semantic.

Three extensibility points cover the expected directions of
future change. Adding a provider requires one new module and one
entry in \texttt{registry.\_PROVIDERS}. Adding an OCR engine
requires one new \texttt{OCREngine} subclass. Adding a schema field
requires one Pydantic attribute plus one entry in
\texttt{FIELD\_NORMALISERS}; the evaluator and reporter iterate
over \texttt{BillExtraction.SCHEMA\_FIELDS}, so they pick up new
fields automatically.

\paragraph{Expected deviations.} Three deviations from this design
are anticipated and will be recorded in Chapter~\ref{chap:impl}~\S\ref{sec:impl-deviations}.
First, the provider-API asymmetry in
\S\ref{sec:design-components-api} is itself a deviation from the
na{\"i}ve design of ``one call shape per provider''; the deviation
is documented here because hiding it would invalidate the
evaluation. Second, some providers may not return
\texttt{latency\_ms} uniformly across SDK versions, and the Timer
wrap in \texttt{src/performance.py} is expected to supersede the
provider field. Third, Tesseract language packs are not
universally installed on arbitrary hosts, and a missing pack
downgrades an OCR run to empty text; this is a ground-truth issue
for that condition, not a bug, and will be handled by the
failure-attribution heuristic.

% ------------------------------------------------------------
\section{Traceability summary}
\label{sec:design-traceability}

Figure~\ref{fig:design-trace} gives the end-to-end chain for each
experimental-requirement class. Every row connects a literature
gap to a requirement, to the design realisation in this chapter,
to the test or evaluation hook that will exercise it. The table is
the single defensible artefact for Trait~5: a marker can verify
the chain without reading the chapters in sequence.

\begin{figure}[ht]
    \centering
    \includegraphics[width=\linewidth]{Chapters/figures/fig-5-7-traceability-matrix.png}
    \caption{End-to-end traceability from literature gap to
    evaluation hook for every experimental requirement.}
    \label{fig:design-trace}
\end{figure}

% ------------------------------------------------------------
\section{Summary and forward bridge}
\label{sec:design-summary}

This chapter turned each requirement of Chapter~\ref{chap:spec}
into a design decision, named an alternative for each decision, and
stated the technical reason for rejection. The single architectural
commitment that dominates the chapter is the modality-isolation
boundary: every other decision exists to make that boundary
auditable by someone other than the author.
Chapter~\ref{chap:impl} takes this design and builds it, recording
the real deviations encountered against the expected deviations in
\S\ref{sec:design-maintenance}. Chapter~\ref{chap:evaluation} tests
whether the design produces the results predicted by the
hypotheses in \S\ref{sec:spec-objectives}, including the
alternative-design comparison that returns to the rejected rows of
Figure~\ref{fig:design-decisions} and evaluates them against the
same data.
.
%  Do NOT add \documentclass or \begin{document} here.
%  Follows Chapter 4 (Specification); precedes Chapter 6 (Implementation).
%  Scope: architecture, design decisions with named alternatives and
%  rejection reasons, pipeline and component design, data model,
%  normalisation, evaluation framework, execution model, maintenance,
%  traceability. Written before implementation; "expected deviations"
%  section flags constraints that will shape Chapter 6.
% ============================================================

\chapter{Design}
\label{chap:design}

% ------------------------------------------------------------
\section{From specification to design}
\label{sec:design-bridge}

Chapter~\ref{chap:spec} fixed what the system must satisfy;
this chapter states how it is realised and what was rejected on the
way. Every design decision below is anchored in a requirement from
\S\ref{sec:spec-fr}--\S\ref{sec:spec-experimental-requirements} and
traces upstream to a gap named in \S\ref{sec:lr-gap}. Four
requirements shape the chapter. ER1 (modality isolation) dictates
that the two pipelines share every component except the one under
test, which constrains the architectural boundary in \S\ref{sec:design-overview}. ER2 (multilingual and prompt-controlled
coverage) dictates that the prompt and the LLM are auditable
artefacts, not inlined strings, which forces the component
abstractions in \S\ref{sec:design-components}. ER6 (null-aware
evaluation) dictates that both predictions and ground truth are
normalised identically before comparison, which drives the
normalisation design in \S\ref{sec:design-normalisation}. ER7
(cross-provider generality) dictates that swapping a provider cannot
change pipeline code, which drives the registry in
\S\ref{sec:design-components}. The remaining decisions fall out from
these four.

% ------------------------------------------------------------
\section{System overview}
\label{sec:design-overview}

Figure~\ref{fig:design-arch} gives the architecture. Three
boundaries matter. \textit{Shared input}
(manifest, loader, PDF, rasteriser) is identical across every
condition, so any difference in raw pixel content is a data-source
effect rather than a pipeline effect. \textit{Modality branch}
contains the single step that varies: either Tesseract followed by
the LLM in text mode, or the LLM in vision mode. \textit{Shared
output} (JSON parse, normalise, validate, persist, evaluate) is
identical across every condition, so any difference in final output
distribution is attributable to the branch content rather than to
the finalisation logic.

The branch is the only place ER1 can be violated, which is why it
is drawn as a single box rather than as two parallel pipelines: two
parallel pipelines have two places to drift. Holding the boundary
tight enough for one provider to be swapped independently of one
modality (\S\ref{sec:design-components}) costs some visual symmetry
in the figure but removes the dominant confound identified in
\S\ref{sec:lr-controlled}.

\begin{figure}[ht]
    \centering
    \includegraphics[width=\linewidth]{Chapters/figures/fig-5-1-architecture.png}
    \caption{End-to-end architecture. Shaded regions are held
    constant across conditions; the central dashed region is the
    modality branch where the input differs. The design isolates
    the ER1 control to a single step rather than to a whole
    pipeline.}
    \label{fig:design-arch}
\end{figure}

Figure~\ref{fig:design-seq} refines the same picture at the
per-document scale. The two paths differ only in whether Tesseract
sits between the rasteriser and the LLM; every call on either side
of that step is identical. The timing keys in brackets match the
per-stage measurements written to \texttt{timings.json} and drive
the latency analysis for RQ2.

\begin{figure}[ht]
    \centering
    \includegraphics[width=\linewidth]{Chapters/figures/fig-5-2-sequence-diagram.png}
    \caption{Per-document sequence for both modalities. The only
    step that varies is whether Tesseract sits in the chain.}
    \label{fig:design-seq}
\end{figure}

% ------------------------------------------------------------
\section{Design decisions and alternatives considered}
\label{sec:design-decisions}

Table~\ref{fig:design-decisions} summarises the ten decisions that
realise the specification, names the alternative considered for
each, and records the specific technical reason for rejection. The
rest of this chapter expands on the decisions that drive the
experiment outcome; the remaining rows are covered by the table.

\begin{figure}[ht]
    \centering
    \includegraphics[width=\linewidth]{Chapters/figures/fig-5-4-decisions-table.png}
    \caption{Design decisions, rejected alternatives and
    traceability to FR, NFR and ER classes.}
    \label{fig:design-decisions}
\end{figure}

\paragraph{Modality-isolation boundary.} The single branch point
keeps ER1 mechanical rather than editorial. Two parallel pipelines
would satisfy the letter of ER1 on day one and violate it silently
on any later change to, say, the prompt, the post-processor or the
JSON parser. The chosen shape makes provenance of any difference
directly inspectable: if the shared input and shared output are the
same byte-for-byte, the difference must come from the branch.

\paragraph{Cross-provider symmetry.} ER7 requires that swapping a
provider cannot alter pipeline code. The \texttt{LLMProvider} ABC
plus registry in \S\ref{sec:design-components} satisfies this
structurally: the pipeline sees only the ABC, so every provider is
exchanged at the level of a config key. The alternative was a
single adapter class with \texttt{if provider == "openai"}
branches; it was rejected because every provider-specific test
would then need to touch pipeline logic, and a future third
provider would require editing the pipeline rather than adding a
file.

\paragraph{Null-aware classification.} ER6 requires the four
classes of the taxonomy defined in \S\ref{sec:spec-evaluation}. The
classification does not live only in the evaluator: it shapes the
failure-diagnosis heuristic in \S\ref{sec:design-evalframework} and
the reporting in Chapter~\ref{chap:evaluation}. The rejected
alternative was a binary exact-match with both-null counted as
correct; it was rejected because on sparse bills (where many
fields are legitimately absent) it can return above 80\% accuracy
on a model that outputs \texttt{null} for every field, which is the
opposite of the discrimination the RQs require.

\paragraph{JSON parse hardening.} Providers do not agree on output
format (\S\ref{sec:lr-eval}). OpenAI Structured Outputs guarantees
a schema-valid JSON object. Anthropic without Structured Outputs
returns text that may be wrapped in markdown fences and may contain
preamble. A retry-with-repair prompt would satisfy this trivially
by asking the model to reformat, but at the cost of changing the
prompt between conditions and breaking ER2. The chosen design is
a four-stage offline fallback (\S\ref{sec:design-execution}) that
never reissues a call.

% ------------------------------------------------------------
\section{Pipeline architecture}
\label{sec:design-pipeline}

\subsection{Shared preprocessing}
\label{sec:design-pipeline-preproc}

Every document enters a single rasteriser before the branch, so the
pixel source for both paths is bit-identical. The alternative would
be to rasterise once per path; it was rejected because PDF
rasterisation is non-deterministic across invocations in the
general case (font hinting, anti-aliasing) and two independent
rasterisations would introduce a second unintended variable.

\subsection{Modality branch}

The OCR path runs Tesseract with the language taken from the
manifest; only the recognised text is forwarded to the LLM. The
vision path sends the first two rasterised pages directly to the
LLM in vision mode. The two-page cap follows from the scope
commitment in \S\ref{sec:spec-scope}: most target fields appear on
the first page or page two. The rejected alternative was to pass
every page; it was rejected because vision tokens scale
near-linearly with page area and the cost would grow quadratically
with document size for no demonstrated accuracy benefit on target
fields.

\subsection{Shared post-processing}

Every LLM response passes through the same four-stage chain: JSON
parse, field-level normalisation, Pydantic validation, per-document
persistence. Parsing and validation are separated so that a parse
failure is distinguishable from a schema violation: the former
keeps the raw output and retries offline
(\S\ref{sec:design-execution}), the latter fails loudly at
evaluation time. Pydantic is chosen over raw \texttt{jsonschema}
because the same model serves as the OpenAI Structured Outputs
schema, the LLM-output container, the normalisation target, the
evaluator input and the reporting envelope; a single Python class
removes five format-translation layers that would otherwise need to
be kept in sync.

% ------------------------------------------------------------
\section{Component design}
\label{sec:design-components}

Figure~\ref{fig:design-strategy} shows the two strategy hierarchies
and the registry that binds them. Every component visible to the
pipeline is an ABC; every concrete implementation is opaque to the
pipeline and instantiated at run time by
\texttt{registry.get\_provider(config)} (lazy import).

\begin{figure}[ht]
    \centering
    \includegraphics[width=\linewidth]{Chapters/figures/fig-5-3-strategy-pattern.png}
    \caption{Provider and OCR strategy hierarchies bound at
    run time by the registry. The rejected alternative
    (hard-coded dispatch) is shown beneath.}
    \label{fig:design-strategy}
\end{figure}

\subsection{LLM provider abstraction}
\label{sec:design-components-llm}

\texttt{LLMProvider} exposes \texttt{extract\_from\_text},
\texttt{extract\_from\_image} and \texttt{get\_model\_id}. Both
extraction methods return a uniform triple
(\texttt{raw\_output}, \texttt{token\_usage}, \texttt{latency\_ms}).
The pipeline calls only the ABC; provider identity enters through
the registry and leaves in the persisted \texttt{model\_id} so that
downstream evaluation can slice by model. The rejected alternative
was a hard-coded \texttt{if}/\texttt{elif} dispatch inside
\texttt{pipeline.py}; its rejection reason is recorded in
Figure~\ref{fig:design-decisions}.

\subsection{OCR engine abstraction}
\label{sec:design-components-ocr}

\texttt{OCREngine} exposes \texttt{extract\_text(image, lang)} and
is implemented by \texttt{TesseractOCR}. Tesseract is the chosen
concrete engine because it executes locally (no confound with a
second paid API in the OCR path) and ships language packs for the
four languages in scope. The abstraction survives the choice: a
later experiment can substitute EasyOCR or PaddleOCR by adding a
module without touching the pipeline, which is what ER5's
failure-attribution story needs to remain falsifiable.

\subsection{Prompt system}
\label{sec:design-components-prompt}

Prompts are versioned files in \texttt{prompts/}, selected by
\texttt{llm.prompt\_file}. The schema description is not written
into the prompt; it is generated from
\texttt{BillExtraction.schema\_description()} and injected through
the \texttt{\{schema\_description\}} placeholder at render time.
The rendered prompt is archived per run (\S\ref{sec:design-execution}) so a diff-only audit proves prompt
identity across conditions, which is what ER2 requires. Inline
string constants were rejected because differences between
conditions cannot be audited by reading a file.

\subsection{Provider-API asymmetry}
\label{sec:design-components-api}

The two providers use different APIs and different mechanisms for
structured output, and the design records the asymmetry rather
than papering over it. OpenAI is called through the Responses API
with Structured Outputs: \texttt{text\_format=BillExtraction}
constrains decoding to schema-valid JSON. Anthropic is called
through the Messages API without Structured Outputs; the same
prompt instructs a JSON response, and the provider's tendency to
wrap JSON in markdown fences is handled by
\texttt{\_strip\_json\_fencing()} before the shared parser runs.

This asymmetry is deliberate. Forcing both providers into
prompt-instructed JSON for the sake of parity would discard the
capability that gives OpenAI guaranteed schema adherence and shift
the risk profile of the experiment silently. The chosen design
preserves each provider's strongest available output mode and
records the raw LLM response per document (NFR8) so that format
failure is separable from content failure in post-hoc analysis. The
evaluation chapter revisits this asymmetry explicitly; a design
choice that hides it would make that analysis impossible.

\subsection{Configuration and CLI}

The runtime is configured by a YAML file (\texttt{config/default.yaml}) with CLI deep-merge overrides
for the handful of flags exercised per run (\texttt{--mode},
\texttt{--provider}, \texttt{--model}, \texttt{--force},
\texttt{--manifest}). The effective configuration after merging is
archived with the run. An environment-variable-only alternative was
rejected because at run-time audit it lacks the single-file
self-description that reproducibility (ER4) assumes.

% ------------------------------------------------------------
\section{Data model}
\label{sec:design-datamodel}

One Pydantic class, \texttt{BillExtraction}
(Table~\ref{fig:schema}), is the contract between prompt, LLM
output, normaliser, evaluator and reporter. Every field is
\texttt{Optional}, which is not a defensive default but the shape
the 4-class null taxonomy assumes: a distinct \texttt{null} is the
only way \texttt{hallucination} and \texttt{omission} can be
separated from \texttt{incorrect}. \texttt{PipelineResult} bundles
the extraction with the artefacts the evaluator and the
cost/latency accounting need (raw LLM output, OCR text, timings,
token usage, model id, pipeline mode); it is not a database row.
\texttt{DocumentEntry} is the resolved manifest row with its file
paths already verified; only entries that pass the six-step manifest
validation (\S\ref{sec:spec-experimental}) reach the pipeline.

Per-layer dataclasses were rejected because every field rename
would then be a five-file change (prompt description, LLM output
parser, normaliser dispatch, evaluator comparison, report column
header) with silent divergence if one file is missed; a single
Pydantic class fails the rename loudly.

% ------------------------------------------------------------
\section{Normalisation}
\label{sec:design-normalisation}

The same normaliser runs over predictions and ground truth before
any comparison. This is the one design choice that most directly
protects the meaning of ``correct'' against surface variation:
\texttt{"\texteuro"} and \texttt{"EUR"}, \texttt{"01/03/2024"} and
\texttt{"2024-03-01"}, \texttt{"OVO Energy Ltd"} and \texttt{"OVO
Energy"} all collapse to the same value before they are compared.
Normalising predictions alone was rejected because it would push
format parity onto the annotators: the set of correct surface forms
would be whatever the normaliser happened to recognise, not the set
of meaningfully correct forms. That trade breaks ER6.

Normalisation dispatches through \texttt{FIELD\_NORMALISERS}, a
mapping from field name to normaliser function. A monolithic
\texttt{normalise(fields)} with a chain of \texttt{if} statements
was rejected because adding a new field would then require editing
the normaliser, the comparator and the reporter in three places;
the dispatch table localises the addition to one entry in one table.
The float parser is the one non-trivial normaliser: it distinguishes
EU and US decimal conventions from the positions of the group and
decimal separators, falling back to standard parsing when both are
absent. LLM-as-judge comparison was rejected as an alternative
because the evaluator becomes non-reproducible the moment it calls
a remote model, breaking ER4.

% ------------------------------------------------------------
\section{Evaluation framework}
\label{sec:design-evalframework}

The evaluation is specified declaratively in
\S\ref{sec:spec-evaluation}; this section states its decision-tree
realisation.

\subsection{Metric set}

Field-level accuracy (over documents), document-level accuracy
(fraction of documents with no incorrect, hallucinated or omitted
field), overall accuracy (all fields across all documents), and
slice accuracy (by language, by utility type) are the four reported
metrics. The split addresses the failure mode named in
\S\ref{sec:lr-eval}: a single aggregate number can hide systematic
language-specific degradation, and coarse-grained metrics that
aggregate across a 12-field schema can hide field-class-specific
degradation (for example, a pipeline that is accurate on
\texttt{provider\_name} but randomly wrong on \texttt{due\_date}).
The six canonical charts listed in \S\ref{sec:spec-evaluation} are
each pinned to a specific research question. Overall-accuracy and
per-field-accuracy answer RQ1 (modality). Per-stage timing and
cost-per-document answer RQ2 (cost and latency).
Accuracy-by-language and accuracy-by-utility type answer RQ3
(multilingual and cross-issuer degradation). No chart exists
without a question.

\subsection{Null-aware classification}

Each field-level comparison emits one of five classifications
(Table~\ref{fig:null-taxonomy}): \texttt{correct},
\texttt{incorrect}, \texttt{hallucination}, \texttt{omission},
\texttt{both-null}. The first four are counted in the accuracy
denominator; \texttt{both-null} is counted as a match but excluded
from the denominator. The exclusion is what prevents a
\texttt{null}-emitting pipeline from scoring highly on sparse
bills. The classification is decided at compare time and persisted
so that the failure-diagnosis heuristic below can operate on it
without re-running comparison logic.

\subsection{Failure diagnosis}
\label{sec:design-evalframework-diagnosis}

For each non-correct field in an OCR-mediated condition, the
heuristic reads the stored \texttt{ocr\_text.txt} and searches it,
case-insensitively, for the ground-truth value. A hit means
Tesseract captured the value and the LLM misused it: the failure
is attributed to the LLM reasoning step. A miss means Tesseract did
not capture the value and no downstream stage could have succeeded:
the failure is attributed to OCR. In vision-mediated conditions no
intermediate text artefact exists, so all non-correct fields are
attributed to the vision model; the asymmetry is recorded and
discussed in Chapter~\ref{chap:evaluation} rather than hidden behind
a weaker heuristic.

Training a classifier for failure attribution was rejected because
it makes the evaluator itself an unvalidated ML component:
disagreements between the classifier and ground truth would then
need their own taxonomy, and an error in the attribution classifier
would be indistinguishable from an error in the pipeline under
test. The substring heuristic is inspectable; its known
false-positive modes (repeated digits on dates) are recorded in
Chapter~\ref{chap:evaluation}.

\subsection{Cross-run comparator}

\texttt{compare\_runs()} reads per-run \texttt{evaluation.json}
files and assembles the accuracy matrix, performance matrix, null
analysis, and slice comparisons for every run under comparison.
The output feeds the factorial analysis specified in
\S\ref{sec:spec-experimental}: the four cells of the $2 \times 2$
matrix are four runs to the comparator, and the comparator is the
only place accuracy across runs is aggregated. No comparison is
ever computed inside the pipeline, so pipeline changes cannot
accidentally rewrite comparison logic.

% ------------------------------------------------------------
\section{Execution model and reproducibility}
\label{sec:design-execution}

Figure~\ref{fig:design-rundir} gives the per-run directory layout.
A run produces one directory, named
\texttt{\{timestamp\}\_\{provider\}\_\{model\}\_\{mode\}}, that
contains the effective configuration, the manifest snapshot at run
time, the exact prompt template, a summary, and a
\texttt{documents/} subdirectory with one folder per document.
Every document folder holds the normalised extraction, the raw LLM
response, the OCR text when applicable, per-stage timings, token
and cost metadata, and on failure an \texttt{error.json} with the
full traceback. The evaluator is a consumer of this directory: it
reads only from disk, never from in-memory objects owned by the
pipeline, which is why pipeline and evaluator are independently
runnable.

\begin{figure}[ht]
    \centering
    \includegraphics[width=0.9\linewidth]{Chapters/figures/fig-5-6-run-directory.png}
    \caption{Per-run directory layout. The evaluator reads only
    from this directory; pipeline and evaluator never share memory,
    so each can be run, replayed and inspected independently.}
    \label{fig:design-rundir}
\end{figure}

Error isolation is mechanical rather than editorial: each document
is processed inside a \texttt{try}/\texttt{except}, and on failure
an \texttt{error.json} with full traceback is written and the run
proceeds. The evaluator skips any document with an
\texttt{error.json}. A crash on document forty therefore loses one
document rather than forty-one.

Figure~\ref{fig:design-parse} gives the JSON parse hardening chain.
The chain is local and offline: each stage rewrites the input
locally and retries the parse, and the final stage raises with a
diagnostic excerpt of the raw output (already archived for the
run). The rejected alternative was a retry-with-repair round-trip
to the model; it was rejected because it mutates the prompt
conditional on provider behaviour, and any prompt mutation breaks
ER2.

\begin{figure}[ht]
    \centering
    \includegraphics[width=0.75\linewidth]{Chapters/figures/fig-5-5-json-parse-fallback.png}
    \caption{JSON parse hardening chain. Every stage is local and
    offline; no stage reissues a call to the model.}
    \label{fig:design-parse}
\end{figure}

Resume is the default: if a document already has an
\texttt{extraction.json}, it is skipped. The \texttt{--force} flag
re-runs everything. This makes partial re-runs cheap, which matters
because API cost per condition is non-trivial (\S\ref{sec:spec-experimental}).

% ------------------------------------------------------------
\section{Maintenance considerations}
\label{sec:design-maintenance}

Five dependency risks are named rather than implied. Poppler
(binary dependency of \texttt{pdf2image}) and Tesseract (with the
four language packs in scope) are system-level dependencies that
must be present on every run host; the \texttt{README} and
\texttt{requirements.txt} list both. Provider SDKs
(\texttt{openai}, \texttt{anthropic}) are pinned in
\texttt{requirements.txt} because both providers make incompatible
API changes between major versions. Pydantic is pinned to v2
because OpenAI Structured Outputs uses v2-specific schema
metadata. The chart backend (\texttt{matplotlib}) is pinned loosely
because its failure mode is visual, not semantic.

Three extensibility points cover the expected directions of
future change. Adding a provider requires one new module and one
entry in \texttt{registry.\_PROVIDERS}. Adding an OCR engine
requires one new \texttt{OCREngine} subclass. Adding a schema field
requires one Pydantic attribute plus one entry in
\texttt{FIELD\_NORMALISERS}; the evaluator and reporter iterate
over \texttt{BillExtraction.SCHEMA\_FIELDS}, so they pick up new
fields automatically.

\paragraph{Expected deviations.} Three deviations from this design
are anticipated and will be recorded in Chapter~\ref{chap:impl}~\S\ref{sec:impl-deviations}.
First, the provider-API asymmetry in
\S\ref{sec:design-components-api} is itself a deviation from the
na{\"i}ve design of ``one call shape per provider''; the deviation
is documented here because hiding it would invalidate the
evaluation. Second, some providers may not return
\texttt{latency\_ms} uniformly across SDK versions, and the Timer
wrap in \texttt{src/performance.py} is expected to supersede the
provider field. Third, Tesseract language packs are not
universally installed on arbitrary hosts, and a missing pack
downgrades an OCR run to empty text; this is a ground-truth issue
for that condition, not a bug, and will be handled by the
failure-attribution heuristic.

% ------------------------------------------------------------
\section{Traceability summary}
\label{sec:design-traceability}

Figure~\ref{fig:design-trace} gives the end-to-end chain for each
experimental-requirement class. Every row connects a literature
gap to a requirement, to the design realisation in this chapter,
to the test or evaluation hook that will exercise it. The table is
the single defensible artefact for Trait~5: a marker can verify
the chain without reading the chapters in sequence.

\begin{figure}[ht]
    \centering
    \includegraphics[width=\linewidth]{Chapters/figures/fig-5-7-traceability-matrix.png}
    \caption{End-to-end traceability from literature gap to
    evaluation hook for every experimental requirement.}
    \label{fig:design-trace}
\end{figure}

% ------------------------------------------------------------
\section{Summary and forward bridge}
\label{sec:design-summary}

This chapter turned each requirement of Chapter~\ref{chap:spec}
into a design decision, named an alternative for each decision, and
stated the technical reason for rejection. The single architectural
commitment that dominates the chapter is the modality-isolation
boundary: every other decision exists to make that boundary
auditable by someone other than the author.
Chapter~\ref{chap:impl} takes this design and builds it, recording
the real deviations encountered against the expected deviations in
\S\ref{sec:design-maintenance}. Chapter~\ref{chap:evaluation} tests
whether the design produces the results predicted by the
hypotheses in \S\ref{sec:spec-objectives}, including the
alternative-design comparison that returns to the rejected rows of
Figure~\ref{fig:design-decisions} and evaluates them against the
same data.
